\documentclass{article}
\usepackage{graphicx} 

\setlength{\parindent}{15pt}
\setlength{\parskip}{5pt}

\usepackage[margin=1.5in]{geometry}
\usepackage{amssymb}
\usepackage{amsmath}
\usepackage{mathtools}
\usepackage{xcolor}
\usepackage{cancel}
\usepackage{centernot}
\usepackage{nicefrac}
\usepackage{bm}
\usepackage{indentfirst}
\usepackage{float}
\usepackage[round]{natbib} 
\usepackage[hidelinks]{hyperref}
\usepackage{enumerate}
\usepackage{titling}
\usepackage{enumitem}
\usepackage{parskip}
\usepackage{subcaption}

\setlist[itemize]{label=\textendash}

\newenvironment{answer}{\par\vspace{0.8em}\textbf{A:}}{\par}
\newcommand{\red}[1]{\textcolor{red}{#1}}
\newcommand{\blue}[1]{\textcolor{blue}{#1}}
\newcommand{\fa}{\text{ for all }} 
\newcommand{\ex}{\text{ there exists }}
\newcommand{\st}{\text{ such that }}
\newcommand{\for}{\text{ for }}
\newcommand{\andt}{\text{ and }}
\newcommand{\ort}{\text{ or }}
\newcommand{\ift}{\text{ if }}
\newcommand{\where}{\text{ where }}
\newcommand{\but}{\text{ but }}
\newcommand{\nat}{\mathbb{N}}
\newcommand{\rat}{\mathbb{Q}}
\newcommand{\com}{\mathbb{C}}
\newcommand{\real}{\mathbb{R}}
\newcommand{\norm}[1]{\lVert#1\rVert}
\newcommand{\abs}[1]{\lvert#1\rvert}
\newcommand{\corr}{\rightrightarrows}
\newcommand{\ul}{\underline}
\newcommand{\bft}{\textbf}
\newcommand{\itt}{\textit}
\newcommand{\E}{\mathbb{E}}
\newcommand{\Prob}{\operatorname{P}}
\newcommand{\Var}{\operatorname{Var}}
\newcommand{\Cov}{\operatorname{Cov}}
\newcommand{\Unif}{\operatorname{Unif}}
\newcommand{\Bias}{\operatorname{Bias}}
\newcommand{\N}{\operatorname{N}}
\newcommand{\MSE}{\operatorname{MSE}}
\newcommand{\iffs}{\Leftrightarrow}
\newcommand{\prefeq}{\trianglerighteq}
\newcommand{\pref}{\triangleright}
\newcommand{\indiff}{\mathrel{\triangleright\triangleleft}}

\title{Conceptual framing}
\date{\today}

\begin{document}

\maketitle

% \begin{center}
%     {\Large Conceptual analysis} \\
%     \vspace{0.5em}
% \end{center}

% \section*{\centering \LARGE \normalfont Codebook}

\section*{Policy choice framework}

Claim: redistributive preferences can be framed as stemming from either beliefs or preferences about the economic game.

Let the \textit{economic game} be defined as the conditional distribution of a society's members' income and wealth depending on their choices of effort, investment, risk-taking, personal traits, and inherited wealth. Knowing this conditional distribution would tell us the current distribution of income and wealth, one's personal circumstances, circumstances of in-groups and out-groups, notions of mobility, and whether competition is zero-sum.

\textit{Beliefs} are disagreements about the current state of the economic game and how the economic game would change in response to policy. 

\textit{Preferences} map a societal economic game to a personal utility. This definition encompasses notions of fairness, generosity, universalism, and preference for consumption vs. leisure. Fairness could be described as the mapping between effort, risk-taking, or investment to income; generosity as the extent to which others' income factors into one's own utility; universalism as the extent to which specific groups factor into one's own utility; and preference for leisure as how much one values economic games that minimize effort. 

\blue{Useful for decomposing age gradient? Possible accounting strategies? }

\subsection*{2x2 beliefs}

Changes in the economic game are often most relevant for young people because they haven't played it yet; they have yet to work hard in their careers, invest in their education, and take risks to advance themselves. 

However, motives for redistribution may depend on which game young people believe themselves to be playing. Consider a simplified situation where the game is defined by whether it is zero-sum vs. positive-sum and high-mobility vs. low-mobility. In a mobile/zero-sum setting,
redistribution might reduce wasteful effort (“rat race” competition), especially for individuals with a strong distaste for effort
or preference for leisure. In immobile settings, redistribution may instead function as insurance, especially for those with high perceived exposure to downside risk.
% insure themselves from unlucky career/income draws

\begin{table}[H]
    \centering
    \begin{tabular}{c|cc}
        & \textbf{Immobile} & \textbf{Mobile} \\
        \hline
        \textbf{Positive-sum} & Allow some risk & Allow some competition \\
        \textbf{Zero-sum} & Insure income & Reduce competition
    \end{tabular}
    \caption{Beliefs about economic game and redistributive preferences}
\end{table}

If exposure to career competition and income risk declines with age, the relative weight on these motives may shift over the
life cycle. That is, even fixing one's underlying beliefs about the system over their entire life cycle, their policy preferences could evolve as they age.


\end{document}
